\scanpage{94}
Les arêtes $\subset t_0$ seront appelées horizontales \deleted{\unsure{basiq.|basiques}}basiques, celles qui lient des él. de $t_0$ à des él. de $S\setminus t_0$ sont appelées verticales, celles qui relient des pts de $S\setminus t_0$ sont appelées horizontales supérieures. \deleted{Il y a \unsure{corresp.|exactement}} L'ens des \deleted{\unsure{Prop. 1|Respectivt|Rappel}} arêtes basiques est $\simeq\unsure{\mathcal P|\mathfrak P}_2(t_0)$, l'ens des arêtes verticales est $\simeq\coprod_{\alpha\in t_0}\widetilde E_\alpha$ (\deleted{\unsure{pour}} à $s\in\widetilde E_\alpha$ correspond l'arête verticale $\{s,\alpha\}$).

Exprimons l'axiome 1 pour les arêtes envisagées jusqu'à présent (celles qui ne sont pas hor. sup.). Pour une arête horizontale inférieure, il est déjà exprimé par le fait que les $\widetilde E_\alpha$ $(\alpha\in t_0)$ sont mutuellement disjoints, i.e. par l'écriture $(\underline{1})$. Pour une flèche verticale $\{\alpha,s\}$ $(\alpha\in t_0,s\in S\setminus t_0)$, celà s'exprime ainsi :

\diagram{\begin{tikzcd}[ampersand replacement=\&,row sep=small,column sep=small]
s\arrow[r,dash]\arrow[dr,dash] \& s'\arrow[d,dash]\\
{} \& a\\
{\cdot\, b} \& {\cdot\, c}
\end{tikzcd}}
\noindent\vrule\hspace{.5em}\begin{minipage}{.89\linewidth}\uline{Prop. 2}\quad $\deleted{\unsure{\forall}}\forall\alpha\in t_0$, et $\forall s\in\overline E_\alpha$, $\exists!\,s'\in\overline E_\alpha$ lié à $\alpha$.\end{minipage}

\noindent\vrule\hspace{.5em}\begin{minipage}{.89\linewidth}\uline{Cor}. \quad $\alpha\longrightarrow\alpha'$ est une involution ss pt fixe de $\widetilde E_\alpha$, notée $\sigma_\alpha$. On dénote par $\sigma$ l'application $S\longrightarrow S$ telle que
\[
\sigma\mid t_0=\mathrm{id}_{t_0},\qquad
\sigma\mid\widetilde E_\alpha=\sigma_\alpha\quad\text{pour }\alpha\in t_0.
\]
\end{minipage}
% archive sheet number: 93
